% Slides 21--40
\sectionframe{02 / QuantLib and Python}{Use a library for conventions; own the economics you extend}{QuantLib v1.43 documentation and source; repository architecture.}

\lesson{QuantLib is a C++ pricing library with language bindings}
{The C++ core supplies instruments, term structures, models, engines, calendars, solvers, and numerical utilities.}
{Python bindings expose much of that object model while Python orchestrates data, workflows, and presentation.}
{The repository runs QuantLib 1.43 through Python, with NumPy, pandas, and Flask around it.}
{The binding boundary is productive when native objects fit; it becomes visible when a bespoke joint process is required.}
{Official QuantLib GitHub release v1.43; requirements.txt.}

\lesson{Instrument, engine, and market data are separate objects}
{An instrument describes contractual cash flows and exercise rights.}
{A pricing engine implements a valuation method and reads market handles or model state.}
{Relinkable handles let market data change without rebuilding every instrument.}
{This separation is the central QuantLib design pattern used throughout the repository.}
{QuantLib architecture; portfolio.py.}

\lesson{The global evaluation date is powerful and dangerous}
{\texttt{ql.Settings.instance().evaluationDate} defines ``today'' for many QuantLib objects.}
{Every curve, schedule, fixing check, and option expiry must agree with it.}
{Tests set it explicitly to prevent machine-date-dependent results.}
{A server must serialize or isolate valuation contexts because this setting is process-global.}
{today.py; portfolio.py valuation\_date; REST integration specification.}

\lesson{Handles express observability and relinking}
{A \texttt{QuoteHandle} wraps a market quote; a \texttt{YieldTermStructureHandle} wraps a curve.}
{Models and engines subscribe to handles and see updated market objects.}
{This supports bump-and-revalue without changing product definitions.}
{It also means hidden shared state can surprise concurrent request processing.}
{QuantLib handles; portfolio.py.}

\lesson{Calendars and schedules deserve native library support}
{Generating coupon dates correctly requires holiday calendars, business-day rules, stubs, and termination conventions.}
{QuantLib's \texttt{Schedule} makes those rules explicit and testable.}
{The repository uses settlement calendars and explicit spot starts rather than date arithmetic alone.}
{Reusing native schedule machinery is safer than recreating conventions in custom Python.}
{portfolio.py \_build\_schedule; standalone\_xccy\_pricer.py \_schedule.}

\lesson{Python is the orchestration layer}
{JSON parsing, schema validation, Flask routes, HTML contexts, diagnostics, and CLI behavior are concise in Python.}
{NumPy handles vectorized path simulation and linear algebra for the custom XCCY engine.}
{pandas structures portfolio tables and curve reports.}
{Python optimizes iteration speed and transparency; it does not erase the need for numerical discipline.}
{Repository source tree and requirements.txt.}

\lesson{C++ is the native extension layer}
{QuantLib models and engines are designed around C++ inheritance, templates, and shared pointers.}
{A production custom stochastic process can live beside existing term structures and Monte Carlo engines.}
{C++ offers tighter control over memory, parallelism, deterministic numerics, and wrapper exposure.}
{The cost is build complexity, ABI management, longer iteration cycles, and a larger validation surface.}
{Official QuantLib source; XCCY pricing specification section 9.}

\twocolesson{Python and C++ solve different bottlenecks}
{Python first}
{Best for proving dynamics, calibration conventions, cash-flow partitioning, LSM policy design, schemas, and diagnostics.}
{C++ next}
{Best when path counts, latency, concurrency, integration policy, or governed library ownership demand it.}
{Prototype the economics in Python; migrate measured bottlenecks and stable interfaces deliberately.}
{standalone\_xccy\_pricer.py; docs/callable\_bermudan\_xccy\_swap\_pricing\_spec.md.}

\lesson{Bindings are not automatic model composition}
{QuantLib exposes \texttt{HullWhiteProcess}, \texttt{BlackScholesMertonProcess}, and \texttt{StochasticProcessArray}.}
{Putting independent processes into an array correlates shocks but does not derive cross-currency drifts for you.}
{The FX drift must be pathwise $r_d-r_f$ and the foreign factor needs its domestic-measure quanto term.}
{A syntactically composable process array can still be economically wrong.}
{Installed QuantLib 1.43 introspection; upstream stochasticprocessarray source.}

\lesson{QuantLib provides multiple swaption engines}
{Analytic engines include Black, Bachelier, Jamshidian, and G2 approaches.}
{Finite-difference engines cover Hull-White and G2 variants.}
{\texttt{TreeSwaptionEngine} handles lattice valuation for a QuantLib \texttt{Swaption}.}
{Engine availability does not imply support for an arbitrary callable two-currency portfolio.}
{Official QuantLib ql/pricingengines/swaption directory, inspected 2026-08-21.}

\lesson{QuantLib now has native constant-notional XCCY instruments}
{Version 1.43 exposes constant-notional cross-currency swaps, basis swaps, and fixed-versus-floating variants.}
{A discounting constant-notional XCCY engine and cross-currency curve helpers are also available.}
{These are valuable for deterministic cash-flow benchmarks and curve construction.}
{They do not provide a Bermudan callable hybrid rate/FX engine.}
{QuantLib v1.43 Python introspection and official crosscurrencyswap source.}

\lesson{A native instrument is not a native callable engine}
{The upstream source contains XCCY instruments and swaption engines as separate families.}
{Code search finds no callable cross-currency swaption engine joining both families.}
{The one-model tree swaption interface targets a swaption on one swap representation.}
{The extension therefore reuses primitives while owning joint dynamics and optimal stopping.}
{Official QuantLib GitHub code search, 2026-08-21.}

\lesson{Repository architecture keeps Flask thin}
{\texttt{app.py} creates the application; the package owns configuration, authentication, routes, and dashboard contexts.}
{\texttt{portfolio.py} remains the core single-currency pricing module.}
{The standalone XCCY engine is separate so its numerical workflow is usable without Flask.}
{Thin delivery layers make model testing and future C++ replacement easier.}
{AGENTS.md; ARCHITECTURE.md; README.md.}

\lesson{JSON makes pricing inputs auditable}
{Versioned market and deal schemas distinguish economics from transport.}
{Unknown fields and unsupported conventions are rejected rather than silently ignored.}
{Normalized input hashing links a result to the exact semantic request.}
{A reproducible pricer records inputs, model version, software versions, seed, and numerical controls.}
{data/schemas; standalone\_xccy\_pricer.py.}

\lesson{Errors should become actionable questions}
{A generic ``bad request'' does not help a trading-system user complete a deal.}
{The MCP validator maps missing or inconsistent fields to precise questions.}
{Statuses separate READY, NEEDS\_INPUT, and INVALID.}
{Explicit incompleteness is safer than guessing a currency, sign, fixing, or exercise rule.}
{xccy\_pricer\_mcp.py; docs/mcp\_xccy\_pricer.md.}

\lesson{Tests should mirror the architecture}
{Curve tests reprice bootstrap helpers and inspect zero/forward outputs.}
{Product tests check schedules, identities, marks, sensitivities, and limiting cases.}
{Flask tests verify authentication, routes, rendering, and state updates.}
{XCCY tests add calibration, martingales, convergence, LSM, and API validation.}
{tests/ directory.}

\lesson{A research workbench is not a production risk system}
{The repository favors transparency, deterministic examples, and inspectable intermediate values.}
{It intentionally omits databases, market feeds, distributed workers, XVA, and enterprise entitlements.}
{Production approval would require independent validation, governed data, observability, access control, and release controls.}
{The documentation states these boundaries instead of implying bank-grade approval.}
{README.md; ARCHITECTURE.md; model limitations in result.json.}

\lesson{Technical architecture follows quantitative ownership}
{QuantLib owns standard conventions, curves, helpers, and benchmark engines.}
{Repository Python owns trade normalization, scenario orchestration, UI, and custom hybrid simulation.}
{Deployment owns process isolation, secrets, TLS, service health, and rollback.}
{Each boundary has different failure modes and therefore different tests.}
{Repository source and deployment templates.}

\lesson{The core engineering rule is explicitness}
{Explicit dates beat implied schedules; explicit signs beat ``payer'' ambiguity across systems.}
{Explicit units beat raw numbers; explicit measures beat drift formulas without numeraires.}
{Explicit tolerances beat a green status without thresholds.}
{Explicit limitations make a research model more useful, not less.}
{AGENTS.md; callable XCCY specifications and schemas.}
